\subsection{Evaluation protocol}
\label{sec:protocol}


Every cohort is evaluated identically, and every decision --- inclusion rule, task definition,
contact rule, unit conversion, split design --- was written into a per-cohort protocol document
\emph{before} any classifier was fitted on that cohort. Where the direction of the study changed
mid-stream, the protocol carries a dated addendum, not a rewrite.

\paragraph{Partitions.} Five repeats of five stratified outer folds, three inner folds for
selection, nested so that representation and hyperparameters are chosen on training participants
only. \textbf{The participant is the unit}: all recordings of a participant are in the same fold,
and participant decisions are the majority vote over the available tasks. Screening and
confirmation use disjoint master seeds; we treat split-to-split variation of $\pm 0.03$ macro F1 as
noise, and report no gain that does not reappear on the confirmation seed. That band was inherited
rather than measured; \cref{sec:results-splitnoise} measures it on twenty fresh seeds per corpus
and reports where it is too narrow.

\paragraph{Statistics, within a cohort.} Confidence intervals are $10{,}000$ label-stratified
participant bootstrap resamples. Paired contrasts use participant-wise swap randomisation.
Multiplicity is corrected by the Holm procedure \citep{holm1979} within families declared in the
protocol.

\paragraph{Spatial controls.} To test whether a gain comes from \emph{where} the signal sits, we
apply two destructive controls at render time, each preserving every per-stroke marginal exactly:
a \emph{circular shift} of the signal values along each stroke, and a \emph{joint permutation} of
the signal values within each stroke. Both leave the ink and the colour histogram identical; only
the correspondence between position and value is destroyed.

\subsection{The acquisition floor}
\label{sec:floor}


\begin{definition}[Acquisition floor]
For a cohort $\mathcal{D}$ and an evaluation protocol $\Pi$, the acquisition floor
$\phi(\mathcal{D}, \Pi)$ is the score that $\Pi$ assigns to a classifier whose only inputs are
properties of the recording process that no behaviour of the subject can produce.
\end{definition}

We instantiate it with two features per recording: the \textbf{median sampling interval between
consecutive contact samples}, and $\log_{10}$ of the \textbf{first timestamp}, that is, the origin
of the session clock. Neither is a symptom: no disease makes a digitiser sample at $9$~ms instead
of $7$~ms, and none sets the clock. The second is independent of the writer by construction, the
first only approximately so, since a digitiser can drop samples under fast motion. That is one
reason to keep the probe to features whose dependence on behaviour is this weak, and to report wider
probes separately.

Either input can be a constant in a given release, and saying which is part of reading the floor:
the median interval varies in one of the six corpora here and the clock origin in three. Where
both are constant the floor is at chance by construction, and the cohort is clean of these two
artefacts by inspection of two columns: a stronger statement than a classifier at chance, not a
weaker one. Where one varies the floor is effectively a one-feature probe: the $0.929$ reported
below on \textsc{Drot\'ar--Dobe\v{s}} is the session clock alone.

The floor is not chance: chance is a property of the label distribution, while the floor is a
property of the dataset as assembled and is the level below which a score is uninformative about the
phenomenon. We therefore report the margin $\Delta_\phi = \text{score} - \phi$ for every arm and
cohort, and treat it rather than the score as the quantity that carries evidence.

Measuring it costs nothing: it reuses the protocol already in place, and the features are two
reductions over a timestamp column. We report it for all six cohorts with a participant endpoint,
including those where our own results are positive.

\subsection{Use of generative AI}
\label{sec:genai}

Generative AI assistance was used in preparing this work and is disclosed here as the journal
requires. A large language model (Anthropic Claude) was used as a coding and editing assistant:
writing and reviewing analysis and rendering code, drafting and copy-editing the manuscript,
checking every bibliographic record against a registry, and auditing the reported numbers against
the run artefacts that produced them. It was not used to design the study, to choose hypotheses or
task subsets, to generate or impute data, or to interpret results; the protocol documents that fix
every design decision predate it and are dated. All code was executed by the authors, all outputs
were reviewed by the authors, and the authors take full responsibility for the content.
% Confirmed by the corresponding author on 2026-11-24 as matching how the tool was actually used
% across the project. Re-check before submission only if another tool is brought in, or if the
% division of labour changes during revision.


